% Ejercicio 2
Para cargar los ficheros de audio se usa la función \texttt{readWave()}:
\begin{lstlisting}[caption={Cargar los ficheros de sonido en R}, label={lst:cargar_ficheros}]
# Audio con el nombre
nombre <- readWave('nombre.wav')
nombre

# Audio con los apellidos
apellidos  <- readWave('apellidos.wav')
apellidos
\end{lstlisting}

Llamar al fichero de sonido (líneas 3 y 7) hace que se muestre información como el número de muestras, la duración y el número de canales. Observar Figura \ref{fig:img2_1}.\\

\begin{figure}[H]
    \centering
    \includegraphics[width=0.95\linewidth]{images/2_cargar.png}
    \caption{Cargar y leer los ficheros de sonido}
    \label{fig:img2_1}
\end{figure}

Los datos obtenidos son los siguientes (ver Cuadro \ref{tab:tabla1}):
\begin{table}[ht]
    \centering
    \caption{Información de los objetos Wave: \texttt{nombre} y \texttt{apellidos}}
    \label{tab:tabla1}
    \begin{tabular}{lll}
        \toprule
        \textbf{Parámetro} & \textbf{Objeto: \texttt{nombre}} & \textbf{Objeto: \texttt{apellidos}} \\
        \midrule
        Número de muestras & 32879 & 49007 \\
        Duración (segundos) & 0.68 & 1.02 \\
        Tasa de muestreo (Hz) & 48000 & 48000 \\
        Canales & Stereo & Stereo \\
        Formato PCM & TRUE & TRUE \\
        Bit & 16 & 16 \\
        \bottomrule
    \end{tabular}
\end{table}

Para obtener el dibujo de la onda la función \texttt{plot} necesita saber el número de muestras que tiene el fichero. Observamos que el fichero \textit{nombre} tiene $32879$ muestras y el fichero \textit{apellidos} tiene $49007$. 

\begin{lstlisting}[caption={Comandos para dibujar las ondas}, label={lst:fun_plot1}]
# Mostrar la onda del sonido
plot( extractWave(nombre, from = 1, to = 32879) )

plot( extractWave(apellidos, from = 1, to = 49007) )
\end{lstlisting}

Al usar la función \texttt{plot} obtenemos la gráfica que se muestra en la Figura \ref{fig:img2_2}:
\begin{figure}[H]
    \centering
    \includegraphics[width=0.95\linewidth]{images/2_2onda.png}
    \caption{Cargar y leer los ficheros de sonido}
    \label{fig:img2_2}
\end{figure}

El dibujo de la onda del archivo de sonido ``nombre.wav'' es el que se muestra en la Figura \ref{fig:img2_3}. El dibujo obtenido para el sonido de ``apellidos.wav'' se observa en la Figura \ref{fig:img2_4}.

\begin{figure}[H]
    \centering
    \includegraphics[width=0.95\linewidth]{images/2_3plot_nombre.png}
    \caption{Dibujo de fichero ``nombre.wav''}
    \label{fig:img2_3}
\end{figure}

\begin{figure}[H]
    \centering
    \includegraphics[width=0.95\linewidth]{images/2_4plot_apellidos.png}
    \caption{Dibujo de fichero ``apellidos.wav''}
    \label{fig:img2_4}
\end{figure}
